\section{Related Work}
\label{sec:related}

\paragraph{Benchmarks for multimodal agents}
Testing digital interaction agents mainly spans coding environments, web scenarios, and mobile applications. In the coding domain, several works provide frameworks and datasets for evaluating agents across programming languages and software engineering activities~\citep{yang2023intercode, jimenez2023swe, li2024devbench, si2024design2code}. 
For web browsing, platforms have been developed for agents to interact with web interfaces through keyboard and mouse actions, alongside datasets focusing on open-ended web tasks and realistic web navigation~\citep{shi2017world, liu2018reinforcement, yao2022webshop, deng2023mind2web, zhou2023webarena, Koh2024VisualWebArenaEM, drouin2024workarena}.
Mobile device interaction research aims at improving accessibility, with simulators for mobile UI interactions and platforms dedicated to InfoUI tasks~\citep{li2020mapping, sun2022meta, venkatesh2022ugif, toyama2021androidenv, rawles2023android, zhang2023mobile, wen2023empowering, zhang2023appagent, wang2024mobile}. 
Further, environments connecting to real computers and datasets for GUI grounding, albeit without interactive capability, have emerged~\citep{gao2023assistgui, Cheng2024SeeClickHG, niu2024screenagent, Kapoor2024OmniACTAD, tan2024towards}. 
Comprehensive task evaluation across different aspects also sees innovations~\citep{liu2023agentbench, mialon2023gaia}.
Differing from previous endeavors focusing on singular environments or lacking executability, \ours integrates an interactive setup enabling agents to engage with operating systems openly, supported by a diverse array of tasks and precise evaluation scripts within a fully controllable setting, marking it as a competitive benchmarking realism and reliability, as well as an environment for learning and evaluating general-purpose digital agent (See Tab.~\ref{tab:benchmark_comparsion} for comparison).


\paragraph{Vision-language models for multimodal agents}
Many existing works on GUI interaction utilize some form of structured data (such as HTML, accessibility trees, view hierarchies) as a grounding source~\citep{deng2023mind2web, gur2023real, li2020mapping, nakano2021webgpt, zhao2022tie, sridhar2023hierarchical, zhang2024large, zhou2023webarena}. 
However, source code often tends to be verbose, non-intuitive, and filled with noise. 
In many cases, it is even inaccessible or unavailable for use, making multi-modality or even vision-only perception a must.
To take screenshots as input, there are already specialized, optimized multi-modal models available that are suited for tasks on web~\citep{Baechler2024ScreenAIAV, furuta2023multimodal, humphreys2022data, lee2023pix2struct, shaw2023pixels} and mobile devices~\citep{hong2023cogagent, zhang2023you}. 
Additionally, general-purpose foundation models~\citep{bai2023qwen, li2023silkie, liu2023visual, zhu2023minigpt} also demonstrate significant potential for multi-modal digital agents.
The development of prompt-based methods~\citep{gao2023assistgui, he2024webvoyager, yan2023gpt, zheng2024gpt}, as well as visual reasoning paradigms, have also further facilitated the performance of digital agents in web pages, mobile apps, and desktop. 
To investigate how well do current models and methods perform in digital agent tasks, our paper evaluates the results of text-only, vision-only, and multi-modal input as well as across multiple methods, demonstrating that existing multi-modal models are far from capable computer agents.
Specifically, there is ample room for improvement in long-horizon planning, screenshot details perception, pixel coordinate locating, and world knowledge.

